%(BEGIN_QUESTION)
% Copyright 2006, Tony R. Kuphaldt, released under the Creative Commons Attribution License (v 1.0)
% This means you may do almost anything with this work of mine, so long as you give me proper credit

Svar på følgende fire spørsmål om deadweight testere:


\vskip 10pt {\narrower \noindent \baselineskip5pt

\noindent
(1) Hva er det ved deadweight tester-prinsippet som gjør instrumentet så nøyaktig og repeterbart? Formulert negativt: hva måtte endres for å påvirke nøyaktigheten i utgangstrykket?

\par} \vskip 10pt




\vskip 10pt {\narrower \noindent \baselineskip5pt

\noindent
(2) Hvorfor er det viktig at en deadweight tester står {\it i vater} under kalibrering av et trykkinstrument?

\par} \vskip 10pt



\vskip 10pt {\narrower \noindent \baselineskip5pt

\noindent
(3) Hvilken effekt har innestengt luft inne i en deadweight tester?

\par} \vskip 10pt




\vskip 10pt {\narrower \noindent \baselineskip5pt

\noindent
(4) Hvorfor er det lurt å rotere primærstempel og vekter {\it forsiktig} mens stempelet er løftet av oljetrykket?

\par} \vskip 10pt



\underbar{file i00154}
%(END_QUESTION)





%(BEGIN_ANSWER)

\vskip 10pt {\narrower \noindent \baselineskip5pt

\noindent
(1) Nøyaktigheten i en deadweight tester bestemmes av tre grunnvariabler. Alle er svært stabile, to kan produseres med høy presisjon, og den tredje er en naturkonstant:

\begin{itemize}
\item{} Massen til kalibreringsvektene
\item{} Arealet til primærstempelet
\item{} Jordens gravitasjon
\end{itemize}

\par} \vskip 10pt




\vskip 10pt {\narrower \noindent \baselineskip5pt

\noindent
(2) Hvis testeren ikke står i vater, vil kraften fra presisjonsvektene ikke være parallell med primærstempelets bevegelsesakse, og stempelet bærer dermed ikke hele vekten korrekt.

\par} \vskip 10pt




\vskip 10pt {\narrower \noindent \baselineskip5pt

\noindent
(3) Innesluttet luft gjør stempelbevegelsen «fjærende» i stedet for fast og stabil.

\par} \vskip 10pt




\vskip 10pt {\narrower \noindent \baselineskip5pt

\noindent
(4) Rotasjon av primærstempelet reduserer virkningen av statisk friksjon, slik at bare dynamisk friksjon (som er mye lavere) påvirker gravitasjonskraften på stempelet.

\par} \vskip 10pt



%(END_ANSWER)





%(BEGIN_NOTES)


\vskip 10pt {\narrower \noindent \baselineskip5pt

\noindent
(2) Hvis testeren ikke står i vater, vil kraften fra vektene ikke være parallell med stempelets akse. Stempelet «ser» da lavere effektiv kraft enn det skal. Tenk på et objekt i en skråning: kraftkomponenten normal på flaten er {\it mindre} enn objektets vertikale tyngde. For deadweight tester gir dette for lavt trykk.

\vskip 10pt

\noindent
I tillegg oppstår sidekraft på føringen til stempelstangen, som skaper friksjon og kan redusere generert trykk ytterligere.

\par} \vskip 10pt





\vskip 10pt {\narrower \noindent \baselineskip5pt

\noindent
(4) Det bør nevnes at {\it forsiktig} rotasjon er alt som trengs for å overvinne statisk friksjon. Noen studenter har en tendens til å spinne vektene som et pottemakerhjul.

\par} \vskip 10pt

%INDEX% Calibration, deadweight tester: basic operation

%(END_NOTES)
